% ============================================================
% COMPARISON NOTE — Paper B (Weyl obstruction spectrum) vs.
%   Abramsky–Cercelescu–Constantin, "Commutation Groups and
%   State-Independent Contextuality," FSCD 2024.
% Purpose: make the "recast" precise -- an explicit equivalence,
% not a descriptive assertion -- and delimit exactly what is new.
% Self-contained; manual thebibliography.
% ============================================================
\documentclass[11pt]{article}
\usepackage[margin=1in]{geometry}
\usepackage{amsmath,amssymb,amsthm,mathtools}
\usepackage[hidelinks]{hyperref}
\usepackage{lmodern}
\newtheorem{theorem}{Theorem}
\newtheorem{proposition}[theorem]{Proposition}
\newtheorem{lemma}[theorem]{Lemma}
\newcommand{\Z}{\mathbb{Z}}
\newcommand{\ip}[2]{\langle #1,#2\rangle}
\title{Paper B and Abramsky--Cercelescu--Constantin:\\ an explicit equivalence, and what is new}
\author{Manuel Flores Gordillo\thanks{ORCID \texttt{0009-0006-8246-0343};
\texttt{manuel.flores@columbia.edu}. Machine check: \texttt{verification/acc\_correspondence.py}.}}
\date{July 2026}
\begin{document}
\maketitle

\begin{abstract}
The value classification of state-independent qudit contextuality --- the even/odd dichotomy and the
fixed obstruction value $d/2$ --- is due to Abramsky, Cercelescu and Constantin~\cite{ACC24}. Paper B
of this program re-derives it in the Weyl--Heisenberg realization. This note makes that relationship
\emph{precise}: we give an object dictionary, prove that the Weyl commutator carry is their commutation
factor under it, and conclude that Paper B's Obstruction Spectrum Theorem is the Weyl specialization of
their Theorem~16 --- neither independent of nor a strengthening of it. We then delimit exactly the
content that is \emph{not} in~\cite{ACC24}. Every correspondence claimed here is machine-checked in
\texttt{verification/acc\_correspondence.py} (PM, and certificates at $d=2,8,16$).
\end{abstract}

\section{The two frameworks}
\textbf{Abramsky--Cercelescu--Constantin~\cite{ACC24}.} A \emph{commutation group} $\mathsf{G}(\mu)$ is
presented by generators $x_1,\dots,x_n$ over $\Z_d$ with a commutator matrix $\mu(x,y)=v_{x,y}$ and a
central phase group $\Z_d$. A \emph{contextual word} $(w,\beta,k)$ is a bracketed word of
commuting subproducts whose evaluated central phase is $k\neq0$; it witnesses state-independent
contextuality (their Thm.~9: one exists iff $\mathsf{G}(\mu)$ is state-independently contextual). Their
\emph{inversion sum} $\sum\mathcal I(w)$ records the central phase accrued in reducing $w$ to normal
form, and (their Lemma~15) for a word of commuting products
\begin{equation}\label{eq:acc}
2\sum\mathcal I(w)\;=\;\sum_{x<y} n_x\,n_y\,v_{y,x},
\end{equation}
$n_x$ the multiplicity of generator $x$. Their Theorem~16 concludes: for a contextual word,
$2k\equiv0\pmod d$ with $k=\sum\mathcal I(w)$, so $d$ must be even and the only nonzero value is $d/2$;
their Theorem~18 gives noncontextuality for odd $d$; Theorem~21 classifies the Darboux-normal-form
matrices that admit contextual words.

\medskip
\noindent\textbf{Paper B (Weyl realization).} Observables are symmetric Weyl operators
$W(v)=\tau^{-q(v)}X^{a_1}Z^{b_1}\otimes\cdots$, $\tau=e^{i\pi/d}$, with symplectic form
$\ip{u}{v}=\sum_i(u_{a,i}v_{b,i}-u_{b,i}v_{a,i})$. A \emph{context} $C$ is a commuting, sum-zero word;
its product is $\omega^{s(C)}I$, $\omega=e^{2\pi i/d}$. A \emph{certificate} $\lambda$
($\lambda^{\mathsf T}A\equiv0$) has value $S=\lambda^{\mathsf T}s\bmod d$, and the \emph{commutator
carry} is $b(C)=\sum_{i<j}\ip{c_i}{c_j}_{\Z}/d\bmod2$. Theorem~1 (Obstruction Spectrum): $2S\equiv0
\pmod d$, so $S\in\{0,d/2\}$ (even $d$), $\{0\}$ (odd $d$).

\section{The dictionary}
\begin{center}
\begin{tabular}{ll}
\hline
\emph{Paper B (Weyl)} & \emph{ACC~\cite{ACC24}}\\
\hline
Weyl label $v$ & commutation-group generator/vector\\
symplectic form $\ip{u}{v}$ & commutator matrix entry $\mu(x,y)=v_{x,y}$\\
closed context $C$ & commuting-product (sub)word\\
certificate $\lambda$ & contextual word $(w,\beta,k)$\\
obstruction value $S$ & central phase $k$\\
commutator carry $b(C)$ & parity of the commutation factor $\sum\mathcal I(w)$\\
$\Z_d$ value assignment & left splitting / noncontextual homomorphism\\
\hline
\end{tabular}
\end{center}

\begin{proposition}[Carry $=$ commutation factor]\label{prop:carry}
Under the dictionary, the Weyl integer pair-sum of a context equals the right-hand side of the ACC
identity~\eqref{eq:acc}:
\[
\sum_{i<j}\ip{c_i}{c_j}_{\Z}\;=\;\sum_{x<y}n_x n_y\,\ip{x}{y}_{\Z}\;=\;2\sum\mathcal I(w).
\]
Hence $b(C)$ is the mod-$2$ reduction of $\sum\mathcal I(w)$, and the Weyl polarization step in the
proof of Theorem~1 is exactly ACC's Lemma~15.
\end{proposition}
\begin{proof}
The first equality is the standard expansion of a sum over an ordered multiset into distinct-generator
multiplicity products (self-pairs contribute $\ip{x}{x}=0$); the second is~\eqref{eq:acc} with
$v_{y,x}=\ip{x}{y}_{\Z}$. Machine-verified on the six Peres--Mermin contexts and on random closed
contexts at $d=2,4,6,8$ in \texttt{acc\_correspondence.py} (all ``[II identical] True'').
\end{proof}

\begin{theorem}[Paper B Thm.~1 is the Weyl specialization of ACC Thm.~16]\label{thm:equiv}
For any Weyl context family, the certificate value $S$ equals the ACC central phase $k$ of the
corresponding contextual word, and Theorem~1's conclusion $2S\equiv0\Rightarrow S\in\{0,d/2\}$ is
Theorem~16's conclusion $2k\equiv0\Rightarrow k\in\{0,d/2\}$. In particular Paper B's Theorem~1 is
neither independent of, nor a strengthening of, ACC's Theorem~16: it is that theorem realized in the
Weyl--Heisenberg group.
\end{theorem}
\begin{proof}
By Proposition~\ref{prop:carry} the two derivations share the polarization identity, and the Weyl
telescoping $2s(C)\equiv-2\sum_{u\in C}q(u)$ paired against a certificate ($\lambda^{\mathsf T}A\equiv0$)
gives $2S\equiv0$, matching $2k\equiv0$; the carried value is $S=k$. The certificate values realized on
concrete AvN witnesses are $S=k=d/2$: Peres--Mermin ($S=1$), the ROZF $d=8$ certificate ($S=4$), and the
fibered $d=16$ certificate ($S=8$), machine-checked in \texttt{acc\_correspondence.py}.
\end{proof}

\section{What is new (and what is not)}
\noindent\textbf{Not new.} The even/odd dichotomy, the value $\{0,d/2\}$, the parity mechanism
$2k\equiv0$, and the classification of which families are contextual are ACC's
(Thms.~16,~18,~21). The carry $b(C)$ is their commutation factor in symplectic-form coordinates: a
change of realization and an explicit closed form, not a new invariant. Their $\Z_2\to\Z_{2k}$
transfer construction is prior art for lifting witnesses up the tower.

\medskip
\noindent\textbf{Candidate new content} (offered at exactly this strength):
\begin{enumerate}
\item an \emph{explicit value-bit closed form} for $2$-adically lifted certificates (Paper B, Thm.~W),
a computation ACC do not perform (they have no $2$-adic lift);
\item the \emph{cross-dimensional} tower statements --- Depth Rigidity and the fiber doubling law ---
which are theorems about relations \emph{between} $d$ and $2d$, and are new only insofar as they are not
corollaries of ``the only nonzero value is $d/2$'' restated at each $d$; the ACC transfer is their
baseline;
\item a three-device superconducting replication of the $d=2$ witness (statistical, assumption-dependent
--- see the repository's hardware notes);
\item the \emph{geometric/cohomological home}: the identification of the obstruction with a canonical
degree-$2$ class of order $n$ on the stack $[X/\mathrm{PU}(n)]$ of maximal abelian subalgebras (companion
paper), which ACC's combinatorial (rewriting-system) treatment does not address. This is the program's
distinctive contribution.
\end{enumerate}

\noindent\emph{Honest status.} Items 1 and 3 are established at the stated (computational / experimental)
strength. Item 2 requires each tower statement to be proved genuinely cross-dimensional; that separation
is the natural next step. Item 4 is the subject of the companion paper and awaits specialist
certification of its Morita / spectral-sequence bookkeeping.

\begin{thebibliography}{9}
\bibitem{ACC24} S.~Abramsky, \c{S}.-I.~Cercelescu, C.~Constantin, \emph{Commutation Groups and
State-Independent Contextuality}, FSCD 2024, LIPIcs \textbf{299}, 28:1--28:20; arXiv:2603.12197.
\bibitem{ORBR} C.~Okay, S.~Roberts, S.~D.~Bartlett, R.~Raussendorf, \emph{Topological proofs of
contextuality}, Quantum Inf.\ Comput.\ \textbf{17}, 1135 (2017).
\bibitem{ROZF} R.~Raussendorf, C.~Okay, M.~Zurel, P.~Feldmann, \emph{The role of cohomology in quantum
computation with magic states}, Quantum \textbf{7}, 979 (2023).
\end{thebibliography}
\end{document}
